% !TEX root = ../algebra_02.tex
\section{Ортогональные матрицы}

\begin{Definition}{Ортогональная матрица}{}
	Матрица $U \in M_n(K)$ называется ортогональной, если $U U^T = E_n$, что равносильно $U^T U = E_n$.
	Множество всех ортогональных матриц обозначается как $O_n(K)$.
\end{Definition}

\begin{Lemma}{Свойство множества ортогональных матриц}{}
	Множество $O_n(K)$ образует подгруппу в полной линейной группе $GL_n(K)$.
\end{Lemma}

\begin{proof}
	Пусть $U_1, U_2 \in O_n(K)$. Проверим замкнутость относительно умножения:
	\[
	(U_1 U_2)(U_1 U_2)^T = U_1 U_2 U_2^T U_1^T = U_1 E_n U_1^T = U_1 U_1^T = E_n
	\] 
	Если $U \in O_n(K)$, то $U^{-1} = U^T$.
	Проверим, что обратная матрица также ортогональна:
	\[
	U^T (U^T)^T = U^T U = E_n
	\] 
	Следовательно, $O_n(K)$ является подгруппой.
\end{proof}

\begin{Proposition}{Свойства строк и столбцов ортогональной матрицы}{}
	Пусть $U = (a_{ij}) \in M_n(K)$. Следующие условия эквивалентны:
	\begin{itemize}
		\item $U$ — ортогональная матрица.
		\item Строки $U[1, :], \dots, U[n, :]$ образуют ОНС в $M_{1 \times n}(K)$.
		\item Столбцы $U[:, 1], \dots, U[:, n]$ образуют ОНС в $K^n$.
	\end{itemize}
\end{Proposition}

\begin{proof}
	Рассмотрим элемент матрицы $U U^T$ на позиции $(k, l)$:
	\[
	(U U^T)[k, l] = \sum_{j=1}^n U[k, j] \cdot U^T[j, l] = \sum_{j=1}^n a_{kj} \cdot a_{lj} = (U[k, :], U[l, :])
	\] 
	Равенство $U U^T = E_n$ выполняется тогда и только тогда, когда скалярное произведение строк равно $1$ при $k=l$ и $0$ при $k \neq l$, то есть когда строки образуют ОНС.
	Аналогичное рассуждение для $U^T U = E_n$ доказывает утверждение для столбцов.
\end{proof}

\begin{Proposition}{Матрица перехода между ОНБ}{}
	Пусть $E$ — ортонормированный базис (ОНБ) в евклидовом пространстве $V$, а $F$ — некоторый базис $V$.
	Пусть $U$ — матрица перехода от $E$ к $F$ ($U = M_{E \to F}$).
	Тогда $F$ является ОНБ тогда и только тогда, когда $U \in O_n(\mathbb{R})$.
\end{Proposition}

\begin{proof}
	Матрица Грама для базиса $F$ выражается через матрицу Грама базиса $E$:
	\[
	\Gamma_F = U^T \Gamma_E U
	\] 
	Поскольку $E$ — ОНБ, то $\Gamma_E = E_n$.
	Следовательно, $\Gamma_F = U^T U$.
	Базис $F$ является ОНБ $\Leftrightarrow \Gamma_F = E_n \Leftrightarrow U^T U = E_n \Leftrightarrow U \in O_n(\mathbb{R})$.
\end{proof}
